\section*{Chapter \ref{ch:b3:interest-groups-and-regulation} --- Interest Groups, Bureaucracy and Regulation}

\begin{solution}{exo:b3:interest-groups-and-regulation:1}
From less than 5~\% to 25~\%: by more than 20 points. % figdata: lic_state, lic_survey, lic_cps
The 2008 survey was commissioned and counted licences from every level of government;
the 2015 value comes from a question added to the official labour-force survey, with
a different sample and wording. The fall from 29 to 22.4~\% may be a change of measure,
not of licensing.
\end{solution}

\begin{solution}{exo:b3:interest-groups-and-regulation:2}
Each spends $3 \times 10/16 = 1.875$ million dalens; together 7.5 million, 75~\% of the rent.
\end{solution}

\begin{solution}{exo:b3:interest-groups-and-regulation:3}
Competitive output 16 at price 4; the monopoly sets $20 - 2Q = 4$, so $Q = 8$ and $p = 12$. The
triangle is $\tfrac12 \times 8 \times 8 = 32$; the rectangle $8 \times 8 = 64$, twice the triangle.
\end{solution}

\begin{solution}{exo:b3:interest-groups-and-regulation:4}
Licensed workers earning no more than similar unlicensed ones; prices no higher
where the licence is stricter; measured quality clearly higher with the licence; or
licences introduced at the request of consumers over the trade's objections.
\end{solution}

\begin{solution}{exo:b3:interest-groups-and-regulation:5}
$Q^* = (100 - 60)/2 = 20$. The bureau sets $100Q - Q^2 = 60Q$, so $Q = 40$ (marginal benefit there is 20, still
positive). The loss on units 20 to 40 is $\tfrac12 \times 40 \times 20 = 400$, equal to the surplus on the first 20,
$\tfrac12 \times 40 \times 20 = 400$: net surplus is zero.
\end{solution}

\begin{solution}{exo:b3:interest-groups-and-regulation:6}
Each firm gains 5 million dalens, each household loses 30: the firms lobby, the
households do not, though their total loss is larger. With 10{,}000 farms each gains
5{,}000, still worth defending, but free riding is harder to overcome without an
association that organises them.
\end{solution}

\begin{solution}{exo:b3:interest-groups-and-regulation:7}
$e^{0.13} - 1 = 13.9~\%$ higher wages. Ability or training not in the controls: if those who
enter licensed occupations would have earned more anyway, the coefficient
overstates the effect of the licence.
\end{solution}

\begin{solution}{exo:b3:interest-groups-and-regulation:8}
Higher demand raises the monopoly price and the profit available at every price;
the regulator lets some of that through and raises the price. By the proposition,
the new support-maximising price is still strictly below the new monopoly price:
consumers keep part of the gain, so the regulator buffers the industry's profit
against the shift.
\end{solution}

\begin{solution}{exo:b3:interest-groups-and-regulation:9}
Fewer injuries, complaints or failures of the work in states with the stricter
licence, among comparable buyers, large enough to be worth the higher prices; or
quality falling where a licence is repealed.
\end{solution}

\begin{solution}{exo:b3:interest-groups-and-regulation:10}
The first-order conditions are $V_1x_2/X^2 = 1$ and $V_2x_1/X^2 = 1$, so $x_1/x_2 = V_1/V_2$ and $x_2 = XV_2/(V_1 + V_2)$.
Substituting, $X = V_1V_2/(V_1 + V_2)$, $x_1 = V_1^2V_2/(V_1 + V_2)^2$ and $x_2 = V_1V_2^2/(V_1 + V_2)^2$. The total is $V_1 \times V_2/(V_1 + V_2) <
V_1$ and $V_2 \times V_1/(V_1 + V_2) < V_2$: less than either stake.
\end{solution}

\begin{solution}{exo:b3:interest-groups-and-regulation:11}
Price equal to average cost: $(12 - Q)Q = 16 + 2Q$, or $Q^2 - 10Q + 16 = 0$, with roots 2 and 8; the cap
is the lower price, $p = 4$ at $Q = 8$. The unregulated monopoly sets $12 - 2Q = 2$: $Q = 5$, $p = 7$, profit
$35 - 10 - 16 = 9$. A cap tied to realised costs lets the firm pass any cost through and
removes its reason to cut costs; a cap set in advance from forecasts lets it keep
what it saves.
\end{solution}

\begin{solution}{exo:b3:interest-groups-and-regulation:12}
The tested solution is \code{ex12_licensing_did.py}, in the chapter's \code{exercises/}
folder. The before--after change in the reforming states is $+0.062$: prices rose with
the common trend. The difference in differences is $-0.064$ (s.e.\ 0.008), close to the
true effect of $-0.06$. % output: ex12_licensing_did
\end{solution}

\begin{solution}{pb:b3:interest-groups-and-regulation:electricians}
% figdata: electricians_stats
\textbf{1.} $100 - 0.05Q = 20 + 0.03Q$: $Q = 1{,}000$ thousand hours, $P = 50$ dalens.

\textbf{2.} $\tfrac12 \times 50 \times 1{,}000 = 25{,}000$ thousand: 25.0 million dalens.

\textbf{3.} $\tfrac12 \times 30 \times 1{,}000$: 15.0 million dalens.

\textbf{4.} Buyers cannot judge an electrician's work before it fails, and bad wiring
endangers neighbours too: a licence could in principle protect them.

\textbf{5.} Raise the quality of the work, and so buyers' willingness to pay, by more
than the triangle and the cost of the training it requires.

\textbf{6.} $100 - 0.05 \times 800 = 60$ dalens an hour.

\textbf{7.} $20 + 0.03 \times 800 = 44$.

\textbf{8.} 16 dalens an hour; $16 \times 800$ thousand $= 12.8$ million dalens.

\textbf{9.} $(60 - 50)/50 = 20~\%$.

\textbf{10.} $\tfrac12 \times 16 \times 200$ thousand $= 1.6$ million dalens.

\textbf{11.} Consumers' surplus falls to $\tfrac12 \times 40 \times 800 = 16.0$ million: a loss of 9.0 million.

\textbf{12.} Producers' surplus is $12.8 + \tfrac12 \times 24 \times 0.8 = 22.4$ million, a gain of 7.4; and $7.4 + 1.6 = 9.0$.

\textbf{13.} The association 7.4 million dalens, the consumers 9.0 million.

\textbf{14.} The association spends $7.4^2 \times 9.0/16.4^2 = 1.83$ million, the consumers $7.4 \times 9.0^2/16.4^2 = 2.23$ million.

\textbf{15.} $7.4/16.4 = 0.451$.

\textbf{16.} $7.4 \times 9.0/16.4 = 4.06$ million dalens.

\textbf{17.} The consumers spend more, having more at stake; the winning probabilities
are in proportion to spending, and spending to stakes, so the side with the smaller
stake wins less often.

\textbf{18.} The consumers' 9 million dalens are spread over all of Brenland's
households, each losing a few dalens; each would free ride on the others'
lobbying, and the contest assumes a consumer group that has already solved that
problem.

\textbf{19.} Any positive spending by the association then wins: the licence passes
with certainty, for almost no rent-seeking spending.

\textbf{20.} $53 - 60 = -7$ dalens an hour.

\textbf{21.} $-7 - (57 - 58) = -6$ dalens an hour.

\textbf{22.} Parallel trends: without the repeal, prices in the repealing regions
would have moved like those in the others.

\textbf{23.} The model predicts a fall from 60 to 50, $-10$; the estimate is $-6$.

\textbf{24.} New electricians take time to train and enter, so the price falls
gradually; or the repeal was partial, keeping some rules.

\textbf{25.} If the licence passes, its social cost is the triangle, 1.6 million dalens,
plus the contest's spending, 4.06 million: 5.66 million dalens a year, against a
transfer of 7.4 million from consumers to electricians.
\end{solution}
